\subsection{Equivariant intersection forms at every coprime prime-adic precision}

\begin{statement}{D.40}{Proved}{Pairing-compatible lifts through the universal abelian coprime covers}
Retain the actual source \(q:\widehat F\twoheadrightarrow D\),
its kernel \(K\), and the specified data
\((\alpha,\lambda,h_i)\) of D.39, and assume that \(D\)
has odd order and \(p\nmid |D|\).
For every \(a\geq1\) put
\[
 E_{p,a}=
 \widehat F/\overline{[K,K]K^{p^a}},\qquad
 E_{p,\mathrm{ab}}=\varprojlim_a E_{p,a}.
 \tag{D.40.1}
\]
There is a continuous peripheral automorphism of
\(E_{p,\mathrm{ab}}\) above \(\alpha\), with the prescribed
multiplier and base conjugators, whose compatible
automorphisms of all \(E_{p,a}\) induce
\(\lambda_p\)-similitudes of the compact cover intersection
form over \(\mathbf Z_p\).
The same maps preserve the full \(D\)-equivariant
group-ring-valued intersection form through
the coefficient automorphism induced by \(\alpha\).
They are compatible with restriction and transfer
on all intermediate covers of this regular cover.

In particular every one of the D.37 weight-five
parameters passes the equivariant intersection-form
test at \(p=2\) and \(p=3\), at every prime-adic
precision.  At \(p=2\) the induced map also preserves
the unique \(D\)-invariant quadratic refinement on
compact homology modulo two.
\end{statement}

\begin{proof}
Use the same integral Cayley chains and boundary
cycles as in D.39.  With coefficients in
\(\mathbf Z_p\), let
\[
 M_p=H_1(K,\mathbf Z_p),\qquad
 B_p=\langle\hbox{oriented boundary cycles}\rangle,\qquad
 C_p=M_p/B_p.
 \tag{D.40.2}
\]
Here \(H_1(K,\mathbf Z_p)\) means the maximal
abelian pro-\(p\) quotient of the actual open
free profinite group.
The Schreier maps identify it with the integral
edge-cycle lattice of rank \(|D|+1\).
The boundary sum of D.39.14 is its only boundary
relation over \(\mathbf Z_p\) as well.
The quotient \(C_p\) is the homology of the
closed oriented compactification of that cover,
so it is free, and its intersection form
\[
                    J:C_p\times C_p\longrightarrow\mathbf Z_p
 \tag{D.40.3}
\]
is perfect alternating and \(D\)-invariant.
Thus \(B_p\) is saturated in \(M_p\).
The averaging operator \(|D|^{-1}\sum_{d\in D}d\)
is defined over \(\mathbf Z_p\); it makes both
the boundary sequence and any chosen projection
onto \(B_p\) \(D\)-equivariantly split.

There is no invariant compact summand:
\[
                         C_p^D=0,\qquad
                         (C_p/p)^D=0.
 \tag{D.40.4}
\]
For the first equality, the invariant part of
the graph sequence has \(M_p^D\) of rank two:
augmentation on \(\mathbf Z_p[D]^D\) is multiplication
by the unit \(|D|\).
The invariant part of the boundary sequence
has \(B_p^D\) of rank \(3-1=2\).
The quotient invariants are exact by averaging
and form a torsion-free quotient \(C_p^D\);
their rank is zero.
The same calculation directly over \(\mathbf F_p\)
proves the second equality.

D.39 gives an \(\alpha\)-semilinear isomorphism
\(L:C_p/p\to C_p/p\).
Both forms
\[
              (v,w)\longmapsto J(Lv,Lw),\qquad
              (v,w)\longmapsto\bar\lambda J(v,w)
 \tag{D.40.5}
\]
are nondegenerate \(D\)-invariant alternating
forms on this same module.
Broto--M\o ller--Oliver, Lemma A.1(b),
applies over every finite field, including
characteristic two: on a self-dual module with
zero invariants for an odd group of coprime
order, such forms are equivariantly isometric
\cite{ACoprimeBMO2012}.
The hypotheses are supplied by
\textup{(D.40.3)--(D.40.4)}.
Choose the resulting actual
\(R\in\operatorname{Aut}_{\mathbf F_p[D]}(C_p/p)\)
and replace \(L\) by \(LR\).  It then satisfies
\[
                    L^{\mathsf t}JL=\bar\lambda J.
 \tag{D.40.6}
\]
This is a finite effective choice: enumerate
the commutant of the two deck-action matrices
and test the displayed matrix equation.
The stated theorem proves termination.

For precision, the absence of an absolutely
self-dual nontrivial constituent can also be
seen directly in every characteristic at issue.
If \(U\) were such a constituent over an
algebraic closure, then
\((U\otimes U)^D\) would be one-dimensional.
The flip acts on it by a nonzero scalar
\(\epsilon\), with \(\epsilon^2=1\).
On the other hand, if \(e_D\) is the averaging
idempotent, then
\[
 \operatorname{tr}(\mathrm{flip}\,e_D)
    =|D|^{-1}\sum_{d\in D}\operatorname{tr}(\rho(d)^2)
    =|D|^{-1}\sum_{d\in D}\operatorname{tr}(\rho(d))
    =0.
 \tag{D.40.7}
\]
Squaring permutes an odd-order group, and
the final average projects onto \(U^D=0\).
This contradicts the trace \(\epsilon\) on
the one-dimensional image of \(e_D\).
Consequently the nontrivial self-dual finite-field
blocks have a nontrivial involution on their
endomorphism fields; the finite-field Hermitian
norm gives their isometry comparison.
The paired non-self-dual blocks have the usual
evaluation pairing.  This explains in particular
why characteristic two introduces no symmetric
nonalternating block in \textup{(D.40.5)}.

We now lift \textup{(D.40.6)} integrally by
explicit equations, without division by two.
Lift its matrix arbitrarily to \(\mathbf Z_p\)
and average:
\[
 T_1=|D|^{-1}\sum_{d\in D}
          \rho(\alpha(d))\,\widetilde L\,\rho(d)^{-1}.
 \tag{D.40.8}
\]
It is exactly \(\alpha\)-semilinear, reduces
to \(L\), and is invertible.
Suppose \(T_a\) already satisfies
\(T_a^{\mathsf t}JT_a=\lambda_pJ\pmod{p^a}\).
The discrepancy
\[
 \Delta_a=
   \frac{\lambda_pJ-T_a^{\mathsf t}JT_a}{p^a}\pmod p
 \tag{D.40.9}
\]
is an alternating \(D\)-invariant matrix.
Let \(U_a\) be its strictly upper triangular
part.  Then \(U_a-U_a^{\mathsf t}=\Delta_a\),
including at \(p=2\).
With \(J_0=\bar\lambda J\bmod p\), set
\[
 H_0=J_0^{-1}U_a,\qquad
 \overline H_a=|D|^{-1}\sum_{d\in D}
                  \rho(d)H_0\rho(d)^{-1}.
 \tag{D.40.10}
\]
The identity
\(H_0^{\mathsf t}J_0+J_0H_0=\Delta_a\)
follows from \(J_0^{\mathsf t}=-J_0\).
Averaging retains that identity because
both \(J_0\) and \(\Delta_a\) are invariant.
Lift \(\overline H_a\) arbitrarily and average
once more over \(\mathbf Z_p\) to obtain a
matrix \(H_a\) commuting exactly with \(D\).
Put
\[
                    T_{a+1}=T_a(1+p^aH_a).
 \tag{D.40.11}
\]
This retains semilinearity and corrects the
form modulo \(p^{a+1}\); the quadratic error
is divisible by \(p^{2a}\), hence by \(p^{a+1}\).
The matrices converge to an invertible
\(\alpha\)-semilinear map \(T_C\) with
\[
                       T_C^{\mathsf t}JT_C=\lambda_pJ
                         \quad\text{over }\mathbf Z_p.
 \tag{D.40.12}
\]

On the boundary lattice use exactly
\[
                 T_B(d b_i)=\lambda_p\alpha(d)h_i b_i.
 \tag{D.40.13}
\]
The actual coset bijection D.39.17 and the
single diagonal relation show that this is
an integral semilinear isomorphism.
Choose a \(D\)-linear section \(C_p\to M_p\)
by averaging any section.
Combine \(T_C\) on that summand with
\(T_B\) on \(B_p\).  This gives an invertible
semilinear \(T_M\) on \(M_p\), retaining all
prescribed boundary frames.
The splitting D.39.11--12 is defined with
the same integral tree chains over
\(\mathbf Z_p\).
Thus
\[
          (m,d)\longmapsto(T_Mm,\alpha(d))
                   \quad\text{on }M_p\rtimes D
 \tag{D.40.14}
\]
is an actual continuous automorphism of
\(E_{p,\mathrm{ab}}\).
Its three peripheral comparisons are the
same equal-power identities as D.39.21;
the norm conjugator D.39.22 is now evaluated
over \(\mathbf Z_p\), where every \(e_i\)
is a unit.  This proves the claimed exact
peripheral identities and their compatibility
in all finite quotients \textup{(D.40.1)}.

To retain the full deck-group information,
define the equivariant intersection form by
\[
 {\mathcal I}(v,w)=
       \sum_{d\in D}J(v,dw)\,d\in\mathbf Z_p[D].
 \tag{D.40.15}
\]
Its coefficient at \(1\) is the original
intersection form.  Semilinearity and
\textup{(D.40.12)} give the exact comparison
\[
 {\mathcal I}(T_Cv,T_Cw)
        =\lambda_p\,\alpha_*{\mathcal I}(v,w),
 \qquad
 \alpha_*\Bigl(\sum a_d d\Bigr)=\sum a_d\alpha(d).
 \tag{D.40.16}
\]
For an intermediate subgroup \(H_C=q^{-1}(C)\),
the covering transfer identifies its open
homology with \(M_p^C\), because \(|C|\)
is a unit.  The transfer of a boundary
circle is the sum of its lifted boundary
circles; projection back multiplies it
by the degree after their actual ramification
indices are included.
All these indices divide \(|C|\) and are units.
Taking the boundary quotient therefore identifies
its compact homology with \(C_p^C\) as well.
The identities
\[
 T_M\sum_{c\in C}c
       =\Bigl(\sum_{c\in C}\alpha(c)\Bigr)T_M,
 \qquad |\alpha(C)|=|C|
 \tag{D.40.17}
\]
prove the stated simultaneous restriction
and transfer comparison, including the
degree factors in the intersection form.

At \(p=2\), average any quadratic refinement
of \(J\bmod2\) over the odd group \(D\).
Its polarization is still \(J\).
Any two invariant refinements differ by
an invariant linear functional on \(C_p/2\),
which vanishes by \textup{(D.40.4)}.
The refinement is thus unique.
Since \(\lambda_2\bmod2=1\), pulling it back
by the semilinear isometry gives another
invariant refinement with the same polarization;
the two are equal.
\end{proof}

The independent program
\texttt{a\_coprime\_intersection.py}
computes the compact cup matrix from every
literal oriented Schreier boundary word.
For two cohomology coordinates \(\chi,\psi\),
it maps a free generator \(j\) to
\((\chi_j,\psi_j,0)\) in the Heisenberg group.
An inverse letter contributes
\(\chi_j\psi_j\), while multiplication adds
the first coordinate times the new second
coordinate.  Each boundary has zero first
two coordinates after passing to compact
cohomology.  Their total central value is
the cup pairing of the filled oriented
surface, by its actual product-of-boundaries
and handle relation.
The saved matrices are alternating,
nondegenerate and invariant under every deck
element; their inverses give the intersection
matrices used in \textup{(D.40.9)--(D.40.11)}.

The program checks the \(C_5\) and \(C_5^2\)
covers, of compact ranks four and twelve,
at both primes two and three.
It constructs all four scalar base
automorphisms with their full boundary maps,
and corrects them through modulus \(2^6=64\)
and \(3^6=729\).
Several cases require a nonidentity
finite-field correction and every subsequent
Hensel correction, so the dyadic step is
actually exercised.
All integral Schreier words, quotient maps,
cup matrices, action matrices, corrections,
and resulting full homology maps are saved
with the exact executed output.
